\chapter{Conclusão}

Este capítulo concluí o estudo sumarisando seus diversos experimentos e resultados alcançados. Ele também dá uma visão de possíveis aplicações dos modelos e técnicas de LVE utilizadas no trabalho realizado. Finalizando, são apresentados possíveis estudos em PCGML que podem ser derivados deste trabalho. 

\section{Discussão dos resultados}

A GAN utilizada em \citep{marioGAN} é pública, assim como o agente A* apresentado em \citet{marioAI} e os níveis utilizados nos conjuntos de treinamento da VLGC \citet{VGLC}. O objetivo deste estudo foi avaliar o impacto da técnica de LVE utilizando três algoritmos bioinspirados diferentes: o CMA-ES, o PSO, e também o NSGA-II para experimentos multi-objetivo conforme \citet{marioGAN} sugere. Além de seguir a sugestão de estudo futuro de \citet{marioGAN} utilizando um algoritmo multi-objetivo na técnica de LVE (no caso o NSGA-II), também foram realizadas comparações dos algoritmos bioinspirados em desempenho e valores atingidos tanto nas funções propostas por \citet{marioGAN} quanto na autoral função-objetivo da equação (4.1) com conjuntos de treinamento de granularidade maior do que em \citep{marioGAN}. Por fim, houve uma preocupação em avaliar jogabilidade e qualidade das soluções, avaliações estas realizadas através de métricas extraidas pelo agente A* de Robin Baumgarten e também por um classificador neural, na tentativa de checar se as soluções geradas seguiam um padrão e se a LVE alteraria este padrão, situação que não se realizou.

Realizando estes diversos experimentos, é possível ilustrar como a técnica de LVE pode ser usada para que os conteúdos do gerador sigam alguma condição imposta, condição esta que se traduz numa função-objetivo. Porém, sem o devido cuidado ou filtro, níveis que não podem ser completados podem ser gerados. Também foram feitos experimentos que fortalecem a hipótese de que o uso da LVE não altera suficientemente o padrão dos conteúdos gerados à ponto de se tornarem fora do padrão do modelo neural, e as evidências visuais levam a acreditar que o modelo gera níveis suficientemente diferentes mesmo utilizando a manipulação de variáveis latentes. 

O uso do modelo explorado treinado com diferentes conjuntos e utilizando funções-objetivos diversas na técnica de LVE gera níveis do jogo \textit{Super Mario Bros} diferentes entre si, e um agente pode ser utilizado para decidir a jogabilidade destes níveis, bem como utilizado dentro da própria LVE se for de execução menos custosa do que o agente A* proposto em \citet{marioAI}.

\section{Possíveis Aplicações}

O modelo proposto pode ser utilizado para gerar níveis jogáveis e interessantes do jogo \textit{Super Mario Bros}, com atributos diversos a depender das funções-objetivos utilizadas na LVE. O modelo, treinado em diversos conjuntos distintos, poderia gerar níveis que seriam filtrados no quesito jogabilidade pelo agente A* de Robin Baumgarten, para então serem apresentados a um jogador humano. Uma aplicação deste tipo tem a possibilidade de tornar um jogo com um fator \textit{repĺay} infinito, em que cada nível seria gerado proceduralmente seguindo alguns padrões humanos estabelecidos previamente no conjunto de treinamento. 

\section{Estudos futuros}

Este estudo considera apenas uma única arquitetura e um único jogo. \citet{marioGAN} sugere que o modelo proposto é genérico o suficiente para generalizar qualquer jogo que possa ser representado no formato proposto pelo repositório aberto da \textit{Video Game Level Corpus}, hipótese esta que poderia ser testada empiricamente. Há também muitas arquiteturas propostas para PCGML, como a CESAGAN proposta em \citet{cesagan} e a MultiGAN proposta em \citet{multigan}, que poderiam ser comparadas com a arquitetura da MarioGAN tanto no jogo \textit{Super Mario Bros} quanto em outros, como o \textit{The Legend Of Zelda} utilizado em \citet{cesagan}.

Além disso, todos os modelos citados para comparação utilizam camadas neurais convolucionais, e não \textit{Vision Transformers}, arquitetura avançada de aprendizado profundo que vem ganhando popularidade nos ultimos anos. \citet{vitgan} propõe o uso de \textit{Vision Transformers} dentro de GANs, mas não na área de PCGML, sendo também uma possibilidade de estudo.

Por fim, já que a VGLC representa os níveis em forma de texto, \textit{Large language models}(LLM) também poderiam ser utilizados para gerar estes níveis. \citet{mariogpt} propõe o uso de LLMs para PCGML utilizando o modelo GPT-2 como base com auxílio de \textit{novelty search} e operadores de mutação customizados, chamando o modelo de MarioGPT. Um estudo poderia ser desenvolvido em cima do modelo MarioGPT, comparando com o modelo explorado neste trabalho (a MarioGAN). Um fato interessante é que \citet{mariogpt} também utiliza o agente A* de Robin Baumgarten para teste de jogabilidade dos níveis, o que já estabelece algo em comum entre os dois estudos.
